\documentclass[10pt,a4paper]{amsart}
\usepackage[margin=19mm]{geometry}
\usepackage{amsmath,amssymb,mathtools}
\usepackage[scr=rsfs,cal=euler]{mathalfa}
\usepackage{xfrac}
\usepackage[unicode,hidelinks,bookmarks=false]{hyperref}
\newcommand{\ie}{i.e.\ }
\newcommand{\cf}{cf.\ }
\newcommand{\resp}{respectively\ }
\newcommand{\Subsection}{\subsection}
\providecommand{\nobreakdash}{\nobreak\hbox{-}}
\setlength{\emergencystretch}{2em}
\multlinegap0pt
\usepackage{mathtools}
\usepackage{amssymb}
\usepackage[shortlabels]{enumitem}
\newtheorem{claim}{Claim}
\newcommand{\F}{\mathcal F}
\title{A solution to the Frankl--Kupavskii conjecture on the Erd\H{o}s--Kleitman matching problem}
\author{Cheng Chi, Yan Wang}
\date{Source: arXiv:2605.06389v3 (CC BY 4.0)}
\begin{document}
\maketitle

\noindent\emph{Source and licence.} A solution to the Frankl--Kupavskii conjecture on the Erd\H{o}s--Kleitman matching problem. Version arXiv:2605.06389v3 (CC BY 4.0), distributed by arXiv under the Creative Commons Attribution 4.0 International licence, http://creativecommons.org/licenses/by/4.0/, as recorded in the arXiv licence metadata for this exact version and shown on the abstract page https://arxiv.org/abs/2605.06389v3. Full licence notice: \texttt{licenses/arxiv-2605.06389-CC-BY-4.0.txt}.\par\smallskip

\noindent\textit{Editorial scope.} Counted unit: the article's claim eq:counting-from-missing (source0.tex lines 808--819). The article's complete original proof of it is delivered with it (source0.tex lines 820--958, one \texttt{proof} environment of 139 lines). No mathematics is written, repaired or re-derived: the statement and the proof are the article's own lines, in the article's own notation, and no retained line is substituted or re-flowed.  No further source-local context is needed: every cross-reference of every retained line resolves inside the two delivered spans.  Nothing was omitted from any retained span, so the package carries no omission record.  No retained line contains a citation, so the wrapper carries no bibliography and no bibliography entry.  No retained line contains a figure environment, an image-inclusion command or drawing code, so no image is delivered and the package declares no asset dependency.  Editorial formatting only, in addition: an AMS \texttt{amsart} preamble that restates the article's own declarations and macros used by the retained lines, namely the environments \texttt{claim} on separate counters, together with the standard packages those lines need.  The retained environments print in this document as Claim 1 in order of appearance, whereas the article prints the same items with the numbering of its own full document.  The complete native source of the article as distributed in the arXiv e-print for this exact version is bundled unchanged as \texttt{sources/200016/source0.tex}.
\begin{claim}\label[claim]{cl:counting-from-missing}
    We have
    \begin{equation}\label{eq:counting-from-missing}
        \sum_{j=0}^{m-1}|B_j(M)|
        \le
        \sum_{j=0}^{m-1}\Lambda_j(a,r)+|M|-\binom rm.
    \end{equation}
    Moreover, if equality holds, then $M=\binom Rm$ and for $0\le j\le m-1$,
    \[
        B_j(M)=\left\{E\in\binom{[n]}j: |E\cap A|\ge m+1-j\right\}.
    \]
\end{claim}

\begin{proof}
    Put \(M^+\coloneqq M\setminus\binom Rm\) and
    \(\xi\coloneqq |M^+|=|M|-\binom rm\).
    For \(E\subseteq[n]\), write
    \(M^+(\overline E)=\{Q\in M^+:Q\cap E=\emptyset\}\).

    We first separate the canonical low sets from the non-canonical ones.  Let \(E\in\binom{[n]}j\) for some $j$ with \(0\le j\le m-1\), and write
    \(u=|E\cap A|\), \(v=|E\cap R|\), so \(j=u+v\).
    If \(E\) is canonical, that is, if \(u\ge m+1-j\), then
    \(r-v=r-j+u\ge m+1+r-2j\).
    Since every \(m\)-set contained in \(R\setminus E\) is missing from \(\F_m\), we get
    \[
        |M(\overline E)|\ge\binom{r-v}{m}
        \ge \binom{m+1+r-2j}{m}.
    \]
    Thus every canonical \(j\)-set belongs to \(B_j(M)\), and the number of
    such sets is \(\Lambda_j(a,r)\).

    Let \(\mathcal X\) be the family of non-canonical low sets which lie in
    one of the \(B_j(M)\).  Explicitly,
    \[
        \mathcal X=
        \bigcup_{j=0}^{m-1}
        \left\{
        E\in B_j(M): |E\cap A|<m+1-j
        \right\}.
    \]
    Therefore
    \(\sum_{j=0}^{m-1}|B_j(M)| = \sum_{j=0}^{m-1}\Lambda_j(a,r)+|\mathcal X|.\)
    It remains to prove \(|\mathcal X|\le\xi\).

    For \(E\in\mathcal X\cap\binom{[n]}j\), define its shortage by
    \(h(E)=m+1-j-|E\cap A|\).  Then \(1\le h(E)\le m+1\), and
    \(r+m+1-2j=r-v+h(E)\).  Since \(E\in B_j(M)\), the extra missing
    \(m\)-sets disjoint from \(E\) satisfy
    \[
        |M^+(\overline E)|
        \ge
        \binom{r-v+h(E)}m-\binom{r-v}m.
    \]
    The shortage $h$ determines how many extra missing $m$-sets are needed
    for $E$ to satisfy the defining inequality for $B_j(M)$.  The minimum possible number
    of such extra missing sets is the following quantity.
    For \(1\le h\le m+1\), define
    \begin{equation}\label{eq:case1-dh}
        D_h(r)=\binom{r-m+1+h}{m}-\binom{r-m+1}{m}=h \frac{r^{m-1}}{(m-1)!}+O_m(r^{m-2}).
    \end{equation}
    The function \(x\mapsto\binom{x+h}{m}-\binom xm\) is increasing for \(x\ge m\).
    Hence every \(E\in\mathcal X\) with shortage \(h\) satisfies
    \(\xi=|M^+|\ge |M^+(\overline E)|\ge D_h(r).\)

    We next compare this required number with the total number of possible
    non-canonical sets of shortage at most $h$.  Let \(N_h\) be the number of
    all non-canonical low sets with shortage at most \(h\).  Explicitly,
    \[
        N_h=
        \sum_{j=0}^{m-1}
        \left|
        \left\{
        E\in\binom{[n]}j:
        1\le m+1-j-|E\cap A|\le h
        \right\}
        \right|.
    \]

    For \(h=1\), the condition \(m+1-j-u=1\) gives \(u=m-j\), and hence
    \(j-u=2j-m\).
    Thus the contribution from $j$-layer is $\binom a{m-j}\binom r{2j-m}$.
    This contribution has degree at most \(j\) in \(s\).
    Therefore the only degree \(m-1\) contribution occurs when \(j=m-1\), where \(u=1\) and \(j-u=m-2\), giving $a\binom r{m-2}$.
    All other layers contribute $o(s^{m-1})$.

    For \(2\le h\le m+1\), the degree \(m-1\) contribution again comes only from \(j=m-1\).
    In that layer, $h(E)=2-|E\cap A|$.
    Thus shortage at most \(h\) implies \(|E\cap A|=0\) or \(|E\cap A|=1\), contributing $\binom r{m-1}+a\binom r{m-2}$ to $N_h$ in total.
    The remaining layers contribute $o(s^{m-1})$.

    Hence, we obtain that
    \[
        N_1=\frac{a r^{m-2}}{(m-2)!}+o(s^{m-1}),
        \text{ and }
        N_h=
        \frac{r^{m-1}}{(m-1)!}
        +
        \frac{a r^{m-2}}{(m-2)!}
        +
        o(s^{m-1})
        \text{ for $2\le h\le m+1$.}
    \]
    By \eqref{eq:case1-dh}, since \(a\le\gamma s\), \(r\ge(m+1-2\gamma)s+1\), and
    \[
        \frac{(m-1)a}{r}
        \le
        \frac{(m-1)\gamma}{m+1-2\gamma}+o(1)
        <1,
    \]
    we have, for $1\le h\le m+1$,
    \[
        D_h(r)-N_h
        \ge
        \left(1-\frac{(m-1)\gamma}{m+1-2\gamma}-o(1)\right)
        \frac{r^{m-1}}{(m-1)!}.
    \]
    By the choice of \(s_0\), for $1\le h\le m+1$, we have
    \(N_h\le D_h(r)-1.\)

    The numbers \(D_1(r),\ldots,D_{m+1}(r)\) are strictly increasing, since
    \(D_{h+1}(r)-D_h(r)=\binom{r-m+h+1}{m-1}>0\).  We compare \(\xi\) with
    this sequence.
    If \(\xi<D_1(r)\), then no non-canonical set can belong to any \(B_j(M)\), so \(|\mathcal X|=0\).
    If \(D_h(r)\le\xi<D_{h+1}(r)\) for some \(1\le h<m+1\), then every set in \(\mathcal X\) has shortage at most \(h\), so
    \(|\mathcal X|\le N_h\le D_h(r)-1\le\xi.\)
    Finally, if \(\xi\ge D_{m+1}(r)\), then \(|\mathcal X|\le N_{m+1}\le D_{m+1}(r)-1\le\xi\).
    Thus \(|\mathcal X|\le\xi\), and hence
    \[
        \sum_{j=0}^{m-1}|B_j(M)|
        \le
        \sum_{j=0}^{m-1}\Lambda_j(a,r)+|M|-\binom rm.
    \]
    This proves \eqref{eq:counting-from-missing}.

    It remains to discuss equality.
    If \(\xi>0\), the preceding argument gives \(|\mathcal X|<\xi\).
    Hence equality in \eqref{eq:counting-from-missing} is impossible unless \(\xi=0\), that is, unless \(M=\binom Rm\).

    Assume now that \(M=\binom Rm\).
    For \(E\in\binom{[n]}j\), again write \(u=|E\cap A|\) and \(v=|E\cap R|\).
    Then \(|M(\overline E)|=\binom{r-v}{m}\).
    Hence \(E\in B_j(M)\) if and only if \(\binom{r-v}{m}\ge\binom{r+m+1-2j}{m}\).
    For \(s_0\) sufficiently large, this is equivalent to \(r-v\ge r+m+1-2j\), which is equivalent to \(u\ge m+1-j\).
    Therefore
    \[
        B_j(M)=
        \left\{
        E\in\binom{[n]}j: |E\cap A|\ge m+1-j
        \right\},
    \]
    as claimed.
\end{proof}

\end{document}
